% Diagrama de arquitectura de la plataforma CRS CataLyst-5 + FPGA
% Compilar con pdfLaTeX (requiere TikZ).
\documentclass[tikz,border=8pt]{standalone}
\usepackage[utf8]{inputenc}
\usepackage[T1]{fontenc}
\usepackage{tikz}
\usetikzlibrary{arrows.meta, positioning, shapes.geometric}
\begin{document}
\begin{tikzpicture}[
  node distance=1.55cm and 1.2cm,
  block/.style={draw, rounded corners, minimum width=2.7cm, minimum height=1.15cm, align=center, fill=gray!10},
  smallblock/.style={draw, rounded corners, minimum width=2.4cm, minimum height=1.0cm, align=center, fill=blue!7},
  line/.style={-{Latex[length=2.5mm]}, thick},
  returnline/.style={-{Latex[length=2.5mm]}, thick, dashed}
]
\node[block] (pc) {Computadora\API en Python};
\node[smallblock, right=of pc] (usb) {Adaptador\USB--TTL};
\node[smallblock, right=of usb] (fpga) {FPGA\Intel Cyclone V GX};
\node[smallblock, right=of fpga] (max) {Adaptador\MAX3232};
\node[block, right=of max] (ctrl) {Controlador original\CRS C500C};
\node[block, below=1.6cm of ctrl] (robot) {Robot articulado\CRS CataLyst-5 (5 DOF)};

\draw[line] (pc) -- node[above, font=\small, align=center] {Comandos\y trayectoria} (usb);
\draw[line] (usb) -- (fpga);
\draw[line] (fpga) -- node[above, font=\small] {UART} (max);
\draw[line] (max) -- (ctrl);
\draw[line] (ctrl) -- node[right, font=\small] {Control del robot} (robot);

% Canal de retorno: no atraviesa la FPGA.
\draw[returnline] (ctrl.north) -- ++(0,1.0) -| node[pos=0.25, above, font=\small, align=center] {Respuesta del controlador\(TX)} (usb.north);
\draw[returnline] (usb.west) -- node[above, font=\small] {Datos recibidos} (pc.east);

\node[below=0.55cm of fpga, font=\footnotesize, text width=5.0cm, align=center] {La FPGA procesa la comunicación de comandos hacia el controlador. El canal de retorno llega a la computadora por el adaptador USB--TTL, sin pasar por la FPGA.};
\end{tikzpicture}
\end{document}
